\documentclass[11pt]{article}
\usepackage[margin=2.6cm]{geometry}
\usepackage{graphicx}
\usepackage{booktabs}
\usepackage{amsmath}
\usepackage{amssymb}
\usepackage[hidelinks]{hyperref}
\usepackage{xcolor}
\graphicspath{{../figs/}}

\newcommand{\commit}[1]{\texttt{\small #1}}
\newcommand{\Pu}{P_{\mathrm{u}}}

\title{\textbf{The RSC Validation Program, Part II:}\\
Precursors, the Separatrix, the Cusp Map, a Second Architecture,\\
and Production Traces (E7--E11)}
\author{Workspace record --- \texttt{Recoverable-Self-Coding/workspace\_demo}\\
companion to arXiv:2606.24861 v4 (PRE WT10366E, in review)}
\date{28--30 August 2026}

\begin{document}
\maketitle

\begin{abstract}
\noindent
Between 28 and 30 August 2026 the remaining items of the recorded RSC
validation menu were run on the local rig (Qwen3.6-27B, 4-bit, with the
Jacobian-lens workspace instrument) and, for the second-architecture
test, on Llama-3.1-8B served by Ollama with no lens at all. Every
experiment was pre-registered: designs and quantitative predictions were
committed to git before any rig data existed, falsifiers stated in
advance, and every miss is scored and diagnosed rather than absorbed.
Headlines: (i)~fluctuation precursors --- rising, detrended queue
variance on the metastable approach to the fold --- are present in both
analyzable feedback arms and absent in matched-load $\alpha{=}0$
controls; (ii)~the separatrix under graded compound shocks is monotone
with its 50\% crossing where the committed model put it, and the
two-component memory asymmetry (window-only $\ll$ backlog-only $<$
compound) lands exactly as predicted; (iii)~the rig-horizon cusp map is
confirmed at two of four spot-checks, with the other two landing inside
drift-shrunk boundary noise --- and the $\pm 50\%$ host-service
systematic is now a named, guarded-against methodology finding;
(iv)~the collapse \emph{transfers to a second model family, lens-free}:
sub-capacity, feedback-specific runaway at exogenous true utilization
$0.79$ while the control stays clean at $0.84$ --- and the paired
negative shows the serving stack's cost model (prefix caching) is a
control parameter of the transition, the load-side dual of
certification; (v)~the collapse ingredients are present in production
agentic traces: a $3.7\times$ error feedback kernel, cascade
over-dispersion against the iid null, and a service law rising from
3.6\,s to 10.8\,s across the context range.
\end{abstract}

\section{Discipline}
All experiments follow the protocol of the Part-I program (E1--E6):
the design and its quantitative prediction are committed to the
repository \emph{before} the rig runs (commit hashes cited per section);
protocol deviations and amendments are committed before the affected
data are scored; failed attempts are archived, never deleted; and every
prediction miss is either confirmed as a miss or attributed to a
\emph{measured} confound, with the diagnostic committed alongside.

Two cross-cutting methodology findings emerged and are now standing
policy:
\begin{itemize}
  \item \textbf{The host-service systematic.} Wall-clock service on the
  rig varies by $\pm 50\%$ between sessions (18.6\,s to 33\,s for the
  same congested prompt class), moving every utilization boundary by
  $\sim\!\pm 0.15$. E7 closed it with in-run calibration (dimensionless
  targets $\rho$, $\theta$ held in \emph{true} units); E8's slope
  verdicts were robust to it; E9's fixed-margin spot-checks were not.
  Every future wall-clock phase calibrates in-run.
  \item \textbf{Serving-stack porting traps.} Porting the loop to a
  second serving stack surfaced four traps (warm-anchor degeneracy,
  cache-poisoned calibration, chat-template contamination, per-arm
  cache-regime asymmetry), each archived with its diagnosis. One of
  them turned into a physics result (Sec.~\ref{sec:e10}).
\end{itemize}

\section{E7 --- Fluctuation precursors (the designed protocol)}
\textbf{Claim.} Approaching the fold, the order parameter's
fluctuations show rising variance and autocorrelation (critical slowing
down); away from the fold, no trend. The first attempt on archival E4
data was honestly inconclusive (no power; mean-drift confound); this is
the designed replacement.

\textbf{Design} (\commit{b3ba4bf}; calibration amendment
\commit{381f68c}). Presence-vs-absence at matched load: a feedback arm
($\alpha{=}0.8$) against an $\alpha{=}0$ control with an \emph{identical
paired exogenous schedule} --- same arrival times, same items --- both
pre-filled to the congested service law, at true $\rho=0.45$,
$\theta=2.62$ held by live calibration. Confound-proofed statistics:
per-window linear detrend; residual variance and AC1 of the queue
length (primary); exactly unbiased Bernoulli excess variance
(secondary); Kendall-$\tau$ trends on the pre-onset segment; censoring
below 40 tasks. Committed prediction: feedback $\tau_{q\mathrm{var}}
\approx +0.30$, ignition $\sim\!66\%$; control $\tau \approx 0$,
ignition $0/200$.

\textbf{Result} (3 seeds; \commit{e760206}). Feedback ignited 3/3
(one censored-fast, as the model gave $\sim\!10\%$ probability); both
analyzable feedback arms show the precursor:
$\tau_{q\mathrm{var}} = +0.49$ and $+0.59$, with last-third variance
growth $7$--$18\times$ the paired controls. Controls: no ignition in
any of three arms; primary flat in seeds 0 and 1 ($-0.050$, $-0.067$
--- the seed-0 value within $0.002$ of the model's prediction). The
one miss: seed-2's control returned $\tau = +0.83$; the committed
diagnostic attributes it quantitatively to measured utilization drift
($\rho$ ramping $0.39 \to 0.62$ over the hour, the largest of all six
arms, with clean dose-response across seeds), and
Fig.~\ref{fig:prec} makes the underlying nuance visible: the drifted
control's \emph{magnitude} stays near-flat while its rank statistic is
tall --- rank-based $\tau$ is magnitude-blind. Caveat recorded;
magnitude primaries and drift regression are the next protocol
revision.

\begin{figure}[htbp]
\centering
\includegraphics[width=\textwidth]{val_precursors}
\caption{E7. Left, center: sliding detrended queue variance over the
pre-onset segment; feedback (red) rises toward ignition (dotted
line), the matched control (blue) stays flat. Right: Kendall-$\tau$
per arm against the committed model values (dashed).}
\label{fig:prec}
\end{figure}

\section{E8 --- The separatrix under graded compound shocks}
\textbf{Claim.} Bistability implies a basin boundary in the
two-component state (backlog $B$, window junk $j$): a shock of size $f$
(drop fraction $f$ of the backlog \emph{and} clean fraction $f$ of the
window) either crosses it or falls back.

\textbf{Design} (\commit{23e01ad}). Sequential protocol, state carried
between arms exactly as on the rig (the E6 lesson); hold at
$\rho=0.35$, deep in the wedge ($0.55$ was rejected in predictor
design: cures spontaneously re-ignited 77\% of the time). Committed:
$P(\mathrm{EXIT}\,|\,f) = 0.21/0.18/0.31/0.55/0.95$ across
$f = 0/0.25/0.5/0.75/1$, separatrix $f^{*}=0.75$; ordering
window-only $0.15 <$ backlog-only $0.61 <$ compound $0.95$.

\textbf{Result} (\commit{b0e819c}; Fig.~\ref{fig:sep}). Monotone:
PINNED, PINNED, AMBIG (backlog shrank, on the curve's rising edge),
VOID at $0.75$ (the sequential design's known self-cure path,
$\sim\!4\%$ in the model --- the one unscored point), EXIT at $1.0$
($B \to 1$, stayed clean). The two-component asymmetry landed exactly:
window-only decisively PINNED (the backlog regenerates the junk);
backlog-only missed EXIT by one task ($B\;23 \to 0$ shock; the junky
window rebuilt only 6); compound EXIT. Both falsifiers dodged. An
attempt-1 run was lost to a 1\%-battery hibernation mid-phase and is
archived; a battery monitor now guards long runs.

\begin{figure}[htbp]
\centering
\includegraphics[width=0.62\textwidth]{val_separatrix}
\caption{E8. The committed model curve with measured verdicts
($\blacktriangledown$). At $f{=}1$ the vertical spread \emph{is} the
two-component result: window-only (blue square) pins, backlog-only
(amber) nearly exits, the compound exits.}
\label{fig:sep}
\end{figure}

\section{E9 --- The $\alpha\times\theta$ cusp map and its spot-checks}
\textbf{Claim.} Mean-field: bistability exists iff
$\alpha(1+\theta) > 1$, closing on $\alpha^{*} = 1/(1+\theta)$.

\textbf{Map} (\commit{23e01ad}; Fig.~\ref{fig:cusp}). Three scan
generations taught the honest formulation: the wedge is
\emph{horizon-dependent} (metastable escape times exceed any fixed
protocol time deep in the wedge), so the committed map is drawn at the
rig's own horizon (2100\,s) with E6 slope verdicts and bisection
crossings. At that horizon the wedge appears between $\alpha = 0.15$
and $0.20$ nearly $\theta$-independently --- where mean-field
$\alpha^{*}$ sweeps $0.433 \to 0.160$ --- because the content channel
holds the wedge open below the mean-field boundary (the E6
state-resolution lesson, now as a map). Wedge width grows with both
$\alpha$ and $\theta$, the direction cusp geometry predicts; the
$\alpha{=}0$ artifact floor is localized at $\rho \gtrsim 0.94$, far
from the $\alpha \ge 0.2$ wedges.

\textbf{Spot-checks} (\commit{84a74c2} committed IGNITED/CALM/EXIT/
PINNED at 81--87\% modal; result \commit{e678c7b}). Measured:
IGNITED\,/\,IGNITED\,/\,PINNED\,/\,PINNED --- two clean confirmations
(S1, S4, both deep in-region at true coordinates) and two unresolved:
afternoon services ran 24.5--33\,s against the 25.2\,s planning
constant, shrinking S2's designed margin from $0.12$ to $0.07$ and
S3's from $0.20$ to $0.03$ --- inside the boundary's coin-flip band,
where single-replicate verdicts decide nothing. The $\alpha$-motion
claim ($l_{\mathrm{down}}$ moving $0.16 \to 0.75$ under an
$\alpha$-quench $0.8 \to 0.2$) is not contradicted: S3 landed \emph{on}
the moved boundary rather than clear of it. Re-run with live
calibration is the recorded follow-up.

\begin{figure}[htbp]
\centering
\includegraphics[width=\textwidth]{val_cuspmap}
\caption{E9. $l_{\mathrm{up}}$ (ignition) and $l_{\mathrm{down}}$
(drain) vs $\alpha$ at three deadlines; shaded region = the bistable
wedge at the rig horizon; dashed = mean-field $\alpha^{*}$. Stars: the
four rig spots at their \emph{measured} true utilizations.}
\label{fig:cusp}
\end{figure}

\section{E10/E10b --- A second architecture, lens-free}
\label{sec:e10}
\textbf{Claim.} The collapse is not an artifact of one model or of the
lens instrument.

\textbf{Setup.} Llama-3.1-8B-Instruct (4-bit) served by Ollama; no
lens, no certification instrument --- correctness, lateness, and queue
dynamics only. Weights came via the Ollama registry after the
HuggingFace LFS edge throttled to $\sim\!48$\,kB/s. Directional
predictions committed (\commit{76a111d}): P1 discontinuous collapse at
some $l_c \le 1.05$ on the feedback arm; P2 down-branch pinning; P3 no
runaway in the $\alpha{=}0$ control.

\textbf{The scored burst run} (attempt 4, \commit{f80a06f}): P1
\emph{failed} --- no load collapse to $l = 1.05$. The committed service
profile localized why: post-cap exogenous tasks served uncached
($\sim\!6.2$\,s) but repair offspring rode Ollama's prefix cache
(sub-second) --- the cascade's re-entrant work was nearly free, so its
branching stayed subcritical. The error channel still separated the
arms decisively ($\Pu\;0.62$--$0.95$ vs control $0.08$--$0.33$), and
the down-branch showed \emph{content} hysteresis without queue
hysteresis ($\Pu = 0.95$ at $l=0.60$ against $0.62$ on the way up).

\textbf{E10b, sustained pressure} (two-branch prediction committed,
\commit{ce8d242}; result \commit{6754134}; Fig.~\ref{fig:llama}). Both
arms pre-capped --- which removes the cache globally (truncation shifts
the prefix every task; repairs $4.82$\,s $\approx$ exo $4.47$\,s), the
correct correspondence to the always-uncached Qwen rig. Branch~A
landed: the feedback arm ran away in the \emph{first} phase --- $B\;0
\to 26$, $\Pu = 0.84$, at exogenous-only true $\rho = 0.79$, driven to
an effective $\sim\!1.4$ by its own offspring --- while the control
stayed clean through true $\rho = 0.84$, a \emph{higher} true load
($B$ ending $0/0/1$; $\Pu \le 0.47$). Sub-capacity, feedback-specific
collapse on a second model family with no instrument in the loop.

\textbf{The pair's joint finding.} Same model, same feedback law:
collapse in the uncached regime, none in the cache-subsidized regime.
Prefix caching suppresses $\theta_{\mathrm{eff}}$ for re-entrant work,
pushing $\alpha_{\mathrm{eff}}(1+\theta_{\mathrm{rep}})$ below the
wedge condition --- \emph{a load-side gate, the dual of certification
on the content side}. The serving stack's cost model is a control
parameter of the transition; the theory predicts its own mitigation,
and both sides of it were measured.

\begin{figure}[htbp]
\centering
\includegraphics[width=0.7\textwidth]{val_secondmodel}
\caption{E10b. Backlog trajectories: feedback (red) runs away within
$\sim$50 tasks at exogenous true $\rho = 0.79$; the control (blue)
stays bounded across all three phases up to true $\rho = 0.84$
(phase boundaries dotted), showing near-critical congestion
fluctuations but no collapse.}
\label{fig:llama}
\end{figure}

\section{E11 --- The ingredients in production traces}
\textbf{Question} (observational; analysis plan committed before
running, \commit{c1ac8be}, corrected pass \commit{9a6ef19}; both
results committed as-is). Are the rig's collapse ingredients present in
real agentic operation? Data: 16 local Claude Code sessions with
$\ge 50$ tool calls (22{,}792 tool\_use$\to$tool\_result pairs);
aggregates only, no message content read.

\textbf{Results} (Fig.~\ref{fig:prod}).
\begin{itemize}
  \item \emph{The feedback kernel is present:} $P(\text{error} \mid
  \text{previous error}) = 0.078$ vs $P(\text{error} \mid \text{previous
  ok}) = 0.021$ --- a $3.7\times$ conditional amplification
  ($n = 514/22{,}262$ transitions).
  \item \emph{Cascades are over-dispersed} against each session's own
  geometric (iid) null: error runs $\ge 3$: 5 observed vs $\sim\!0.3$
  expected.
  \item \emph{The service law is present:} median LLM step latency
  rises $3.6 \to 10.8$\,s across the context range (top bin
  $n = 11{,}726$).
  \item \emph{Honest null:} within-session error accumulation is flat
  ($\sim\!0.02$ across progression deciles). Consistent with --- not
  proof of --- the theory's \emph{gated} regime: these sessions have a
  human certification gate in the loop, exactly where junk should not
  accumulate. Pass-1 design errors (confounded context binning; tool
  latency measured instead of LLM step latency) are reported and were
  corrected in a separately pre-committed second pass.
\end{itemize}

\begin{figure}[htbp]
\centering
\includegraphics[width=\textwidth]{val_production}
\caption{E11. Left: the production service law. Center: the feedback
kernel. Right: cascade run-length counts against the iid null
(log scale).}
\label{fig:prod}
\end{figure}

\section{Program status and synthesis}
\begin{table}[htbp]
\centering\small
\begin{tabular}{@{}llll@{}}
\toprule
Experiment & Claim & Verdict & Caveats \\
\midrule
E7 precursors & critical slowing down & \textbf{confirmed} (2/2 arms) &
$\tau$ primary 1/2 pairs; drift-attributed \\
E8 separatrix & basin boundary, 2-component & \textbf{confirmed} &
$f{=}0.75$ point VOID (self-cure path) \\
E9 cusp map & wedge structure in $(\alpha,\theta)$ & 2/4 spots clean &
2 spots inside drift-shrunk margins \\
E10b transfer & second architecture & \textbf{confirmed} &
hysteresis (P2) not yet probed on Llama \\
E10 pair & caching = load-side gate & \textbf{measured both sides} &
--- \\
E11 production & ingredients in the wild & 3 of 4 present &
observational; gated-regime null \\
\bottomrule
\end{tabular}
\caption{The Part-II program at a glance. Everything pre-registered;
all data, predictions, diagnostics, and archived failures in the
workspace repository (through commit \commit{2a1ce3e}).}
\end{table}

Together with Part I (E1--E6: the measured store/encoding/commitment
architecture, loop-level collapse with hysteresis and two-component
memory, the zero-fit quantitative closure, emergent-$\alpha$
discrimination of re-consumption policies, and the gating threshold as
state-resolved structure), the framework's load-side physics now has:
the fold, its precursors, its basin boundary, its control surface, a
second-architecture replication without the instrument, a measured
architectural mitigation predicted by the theory itself, and its
ingredient curves in production traces. The recorded open items ---
an E9 spot re-run under live calibration, a Llama hysteresis probe,
E8 seed 1, and independent replication --- are refinements, not gaps in
the chain of evidence.

\end{document}
